\chapter{\texttt{core.SM} --- Standard Model specialisation (3 flavours)}
\label{ch:core-sm}

\section{\texttt{pmns.py} --- the Standard Model mixing matrix}

\modulehead{core/SM/sm\_pmns.py}{\pyclass{PMNS\_SM}: the concrete 3-flavour
implementation of the abstract class \pyclass{core.common.pmns.PMNS}.}

\begin{theorybox}
\subsection*{The Standard Model product structure}
This is the only module that explicitly fixes a product structure for the
generic factors defined in \cref{ch:core-common}. The full 3-flavour PMNS
matrix is
\begin{equation}
  U = R_{23}(\theta_{23})\,\Delta(\delta)\,R_{13}(\theta_{13})\,\Delta(\delta)^\dagger\,R_{12}(\theta_{12}),
  \qquad \Delta(\delta) = \diag(1,1,e^{i\delta}),
\end{equation}
built from the generators $R_{12}$, $R_{13}$, $R_{23}$, $\Delta$ defined in
\cref{ch:core-common}. \pyfunc{PMNS\_SM.outer\_block} and
\pyfunc{PMNS\_SM.reduced} identify
\begin{equation}
  O \equiv R_{23}(\theta_{23})\,\Delta(\delta),
  \qquad
  U_{\rm red} \equiv R_{13}(\theta_{13})\,\Delta(\delta)^\dagger\,R_{12}(\theta_{12}),
  \qquad U = O\,U_{\rm red},
\end{equation}
so that this single factorisation is the concrete instance, for the SM,
of the abstract $U=O\,U_{\rm red}$ contract from \cref{ch:core-common}.

\subsection*{Why the outer block commutes with the matter Hamiltonian}
$O$ does not mix the electron-flavour index: $R_{23}$ acts only on the
$(\mu,\tau)$ block and $\Delta$ is diagonal, so $O$ has the block form
$O = \begin{psmallmatrix}1 & 0\\ 0 & O_{2\times2}\end{psmallmatrix}$ in the
flavour basis. Since the matter potential only has a non-zero $ee$ entry,
$H_{\rm mat}=\diag(V,0,0)$, it follows that
\begin{equation}
  O\,H_{\rm mat}\,O^\dagger = H_{\rm mat} \qquad \text{(they commute exactly).}
\end{equation}
This is precisely the property \cref{ch:core-common} assumes, without
proof, to derive the general factorisation \cref{eq:H-factorization}; here
it is verified explicitly for the concrete SM $O$.

\subsection*{Transforming the Hamiltonian to the reduced basis}
Given the flavour-basis Hamiltonian
$H_{\rm flavour}=H_{\rm kin}+H_{\rm mat}=U\diag(k_1,k_2,k_3)\,U^\dagger + H_{\rm mat}$
(\texttt{hamiltonian.py} section, \cref{ch:core-common}), substituting
$U=O\,U_{\rm red}$ and using $O\,H_{\rm mat}\,O^\dagger=H_{\rm mat}$ gives
\begin{equation}
  H_{\rm flavour} = O\underbrace{\Bigl[U_{\rm red}\diag(k)\,U_{\rm red}^\dagger + H_{\rm mat}\Bigr]}_{\displaystyle H_{\rm red}}O^\dagger = O\,H_{\rm red}\,O^\dagger,
\end{equation}
matching \cref{eq:H-factorization}: \pyfunc{PMNS\_SM.reduced\_basis}
applied to $H_{\rm flavour}$ (inherited from \pyclass{PMNS}, computing
$O^\dagger(\cdot)O$) yields exactly $H_{\rm red}$ above.

\subsection*{The CP phase drops out of the reduced kinetic term}
$R_{12}$ does not mix the $\tau$ index (third index) and $\Delta$ acts
non-trivially only on that same index, so the product
$R_{12}\diag(k)R_{12}^T$ has no non-zero entries connecting to index 3;
conjugation by $\Delta$ therefore acts trivially on it:
\begin{equation}
  \Delta^\dagger\bigl(R_{12}\diag(k)R_{12}^T\bigr)\Delta = R_{12}\diag(k)R_{12}^T,
\end{equation}
so that \textbf{the CP phase $\delta$ drops out completely from the reduced
kinetic term}, which reduces to the real matrix
$U_{\rm red}^{(0)} \equiv R_{13}R_{12}$:
\begin{equation}
  \keyresult{U_{\rm red}\diag(k)\,U_{\rm red}^\dagger = \bigl(R_{13}R_{12}\bigr)\,\diag(k)\,\bigl(R_{13}R_{12}\bigr)^T.}
\end{equation}
This is why \pyfunc{PMNS\_SM.reduced} returns $R_{13}R_{12}$ rather than
the formally-defined $R_{13}\Delta^\dagger R_{12}$: the two coincide once
contracted with the (real, diagonal) kinetic term, but $R_{13}R_{12}$ is
real, cheaper to evolve, and avoids carrying a phase that always cancels.

\subsection*{Recovering the flavour basis}
Once $H_{\rm red}$ (and its evolutor $S_{\rm red}$) has been obtained, the
flavour-basis result is recovered by dressing once with $O$
(\pyfunc{PMNS\_SM.flavour\_basis}, inherited from \pyclass{PMNS}, computing
$O(\cdot)O^\dagger$):
\begin{equation}
  H_{\rm flavour} = O\,H_{\rm red}\,O^\dagger,
  \qquad
  S_{\rm flavour}(x) = O\,S_{\rm red}(x)\,O^\dagger.
\end{equation}
\end{theorybox}

\begin{keyeqbox}
\textbf{Coherent propagation scheme.} In vacuum ($H_{\rm mat}=0$), a
coherent flavour state is most simply understood, and computed, entirely
in the mass basis:
\begin{enumerate}
  \item \textbf{Flavour $\to$ mass}: project onto mass eigenstates,
    $\psi_{\rm mass}(x_0) = U^\dagger\,\psi_{\rm flavour}(x_0)$;
  \item \textbf{Coherent propagation}: each mass eigenstate accumulates its
    own phase independently, since the mass basis diagonalises the kinetic
    term by construction,
    $\psi_{\rm mass}(x) = \diag\bigl(e^{-ik_1x},\dots,e^{-ik_3x}\bigr)\,\psi_{\rm mass}(x_0)$;
  \item \textbf{Mass $\to$ flavour}: project back,
    $\psi_{\rm flavour}(x) = U\,\psi_{\rm mass}(x)$.
\end{enumerate}
Composing the three steps gives the familiar vacuum evolutor
$S_{\rm vacuum}(x)=U\diag\bigl(e^{-ik_ix}\bigr)U^\dagger$
(\texttt{medium.vacuum.evolutor}). In matter, $H_{\rm mat}$ does not
commute with the full $U$, so this three-step picture does not apply
directly at the full-$U$ level; the reduction trick derived above is
exactly what restores an analogous, computationally cheap three-step
scheme -- flavour $\to$ reduced $\to$ propagate $\to$ flavour -- using $O$
and $U_{\rm red}$ in place of $U$.
\end{keyeqbox}

\begin{implbox}
\pyclass{PMNS\_SM(params: PMNSParams)} implements only the three abstract
methods that depend on the product structure: \pyfunc{pmns\_matrix} (the
full $3\times3$ $U$), \pyfunc{reduced} ($U_{\rm red}=R_{13}R_{12}$), and
\pyfunc{outer\_block} ($O=R_{23}\Delta$). Everything else (rotation
builders, \pyfunc{flavour\_basis}/\pyfunc{reduced\_basis}, antineutrino
selection via \pyfunc{select\_antinu}) is inherited unchanged from the base
class, without duplicating code. It is the default scenario for the whole
project: any function in \texttt{medium.*} that receives a generic
\pyclass{PMNS} object works unmodified with either \pyclass{PMNS\_SM} or
\pyclass{PMNS\_sterile} (\cref{ch:core-bsm}), thanks to the inheritance
design of \cref{ch:core-common}.
\end{implbox}

\begin{refbox}
Standard PMNS parametrisation: \cite{ParticleDataGroup}. Reference numerical
values (global best fits) used in the default presets:
\cite{Esteban2020NuFit,NuFit52}.
\end{refbox}

\section{\texttt{mass\_spectrum.py} --- Standard Model mass spectrum}

\modulehead{core/SM/sm\_mass\_spectrum.py}{\pyclass{MassSpectrum\_SM}, the
concrete 3-flavour implementation of \pyclass{MassSpectrum}
(\cref{ch:core-common}).}

\begin{theorybox}
The Standard Model scenario has no sterile state to accommodate, so
\pyfunc{difference\_vector()} simply returns
\pyfunc{difference\_vector\_base()} unchanged, sized to exactly 3
components:
\begin{equation}
  \keyresult{\bigl(\Delta m^2_1,\ \Delta m^2_2,\ \Delta m^2_3\bigr)
  =\begin{cases}
      \bigl(0,\ \Delta m^2_{21},\ \Delta m^2_{3\ell}\bigr) & \Delta m^2_{3\ell} > 0\ \text{(NO)},\\[2pt]
      \bigl(-\Delta m^2_{21},\ 0,\ \Delta m^2_{3\ell}\bigr) & \Delta m^2_{3\ell} < 0\ \text{(IO)}.
    \end{cases}}
\end{equation}
The ordering-dependent construction itself -- selecting between the two
branches above from the sign of $\Delta m^2_{3\ell}$ -- is entirely
inherited from \pyclass{MassSpectrum} (\cref{ch:core-common},
\texttt{mass\_spectrum.py} section): \pyclass{MassSpectrum\_SM} adds no
logic of its own beyond fixing the output size to 3, in contrast with
\pyclass{MassSpectrum\_BSM} (\cref{ch:core-bsm}), which appends a fourth,
sterile entry $\Delta m^2_{41}$ to the same base vector.
\end{theorybox}

\begin{implbox}
\pyclass{MassSpectrum\_SM(DeltamSq21, DeltamSq3l)}: adds no fields beyond
the base \pyclass{MassSpectrum}; \pyfunc{difference\_vector()} is a
single-line override, \code{return
self.difference\_vector\_base(context=context)}.
\end{implbox}
